% =============================================================================
% CAPÍTULO 9: Orquestração de Pipelines de Validação no OpenCode
% Volume 2 — Guia prático de instalação, configuração e execução
%            de pipelines TDD/SDD com o OpenCode Ecosystem
% =============================================================================
\destaque{Nível-alvo N2 — Pesquisa Reprodutível: instale, configure e opere
pipelines completos de validação de Gêmeos Digitais Periodontais usando
agentes MCP, suítes TDD e integração contínua no OpenCode Ecosystem.}
\label{ch:opencode}
% Seções:
%   9.1  - Instalação e configuração do ambiente OpenCode
%   9.2  - Agentes MCP: catálogo e funções no SUS-Twin
%   9.3  - Orquestração de agentes de validação
%   9.4  - Execução de TDD/SDD no pipeline SUS-Twin
%   9.5  - Pipeline CI/CD para validação contínua
%   9.6  - Badge de validação e relatórios
%   9.7  - Exercícios práticos

% =============================================================================
\section{Instalação e Configuração do Ambiente OpenCode}
\label{sec:opencode-instalacao}

O OpenCode Ecosystem\footnote{Ecossistema \textit{open-source} de orquestração
multiagente composto por mais de 600 componentes validados, incluindo agentes
de inteligência artificial, servidores MCP, motores de raciocínio formal e
pipelines de validação acadêmica.} é a plataforma que materializa a
orquestração de todos os pipelines de validação do Framework SUS-Twin. Antes
de executar qualquer validação, é necessário preparar o ambiente computacional.

\subsection{Pré-requisitos de infraestrutura}

A tabela a seguir consolida os requisitos mínimos e recomendados para executar
o ecossistema em cenários de pesquisa reprodutível:

\begin{table}[ht]
\centering
\caption{Pré-requisitos de infraestrutura para o OpenCode Ecosystem}
\label{tab:opencode-requisitos}
\begin{tabular}{LCC}
\toprule
\thead{Componente} & \thead{Mínimo} & \thead{Recomendado} \\
\midrule
Sistema operacional & Windows 10 / Ubuntu 22.04 & Windows 11 / Ubuntu 24.04 \\
Processador & 4 núcleos, 2,5~GHz & 8 núcleos, 3,5~GHz \\
Memória RAM & 8~GB & 32~GB \\
Armazenamento & 20~GB SSD & 100~GB NVMe \\
Python & 3.10 & 3.12 ou 3.13 \\
Node.js & 18 LTS & 22 LTS \\
Git & 2.30 & 2.45 \\
\bottomrule
\end{tabular}
\end{table}

\subsection{Pipeline de instalação}

O OpenCode CLI\footnote{Interface de linha de comando que gerencia a instalação,
atualização e orquestração de todos os componentes do ecossistema.} é o ponto
de entrada único. A instalação segue três etapas:\footnote{Execute cada bloco
em um terminal separado ou sequencialmente, aguardando a conclusão de cada
etapa antes de prosseguir.}

\begin{lstlisting}[language=bash, caption={Instalação completa do OpenCode Ecosystem},
label={lst:opencode-install}, style=bashstyle]
# Etapa 1: Instalar o OpenCode CLI via npm
npm install -g @opencode/cli

# Etapa 2: Verificar a instalação
opencode --version
# Saída esperada: v1.14.x (ou superior)

# Etapa 3: Inicializar o ecossistema no diretório de trabalho
cd projetos/sus-twin-validador
opencode init --ecosystem full

# Saída esperada:
#  [OK] OpenCode Ecosystem inicializado
#  [OK] 46 MCP servers configurados
#  [OK] 227 skills catalogadas
#  [OK] 128 agentes registrados
\end{lstlisting}

\subsection{Arquivo de configuração central: \texttt{opencode.json}}

Todo o comportamento do ecossistema é governado por um único arquivo JSON
na raiz do projeto. O \autoref{lst:opencode-json} mostra a configuração
mínima necessária para o pipeline SUS-Twin:

\begin{lstlisting}[basicstyle=\ttfamily\scriptsize, caption={Configuração mínima do OpenCode para validação SUS-Twin},
label={lst:opencode-json}]
{
  "$schema": "https://opencode.ai/schemas/v1.14/opencode.json",
  "version": "1.14",
  "ecosystem": {
    "name": "sus-twin-validador",
    "mode": "validation"
  },
  "mcp": {
    "servers": [
      "filesystem", "code-runner", "sequential-thinking",
      "websearch",  "sqlite",       "github",
      "pdf",        "scihub",       "eslint"
    ]
  },
  "agents": {
    "autoevolve":  { "enabled": true,  "interval": "daily" },
    "code-reviewer": { "enabled": true, "strictness": "high" },
    "test-engineer": { "enabled": true, "coverage": 90 },
    "quantum-nexus": { "enabled": false }
  },
  "tdd": {
    "suites_path": "./tdd/",
    "auto_validate_on_commit": true,
    "required_passed_ratio": 1.0
  }
}
\end{lstlisting}

\begin{highlight}
\textbf{Dica:} após editar o \texttt{opencode.json}, execute
\texttt{opencode validate --config} para verificar a integridade do arquivo
antes de iniciar os pipelines.
\end{highlight}


% =============================================================================
\section{Agentes MCP: Catálogo e Funções no SUS-Twin}
\label{sec:opencode-mcp}

O \textbf{MCP} (Model Context Protocol)\footnote{Protocolo padronizado que
permite que agentes de inteligência artificial se conectem a ferramentas e
fontes de dados externas de forma segura e estruturada, similar a um
\textit{plug-and-play} de serviços.} é o barramento de integração do
ecossistema. Cada servidor MCP expõe um conjunto de ferramentas que os agentes
consomem sob demanda. No contexto do SUS-Twin, os seguintes MCPs são
essenciais:

\begin{table}[ht]
\centering
\caption{Agentes MCP essenciais e suas funções no pipeline SUS-Twin}
\label{tab:opencode-mcp-sustwin}
\begin{tabular}{LLL}
\toprule
\thead{MCP / Agente} & \thead{Função no SUS-Twin} & \thead{Exemplo de uso} \\
\midrule
\texttt{filesystem} & Gerencia arquivos de configuração
  e resultados & Ler \texttt{objetivos.json} \\
\texttt{code-runner} & Executa scripts Python isolados
  em sandbox\footnotemark & Rodar \texttt{tdd\_validate.py} \\
\texttt{sequential-thinking} & Mantém contexto de raciocínio
  longitudinal do paciente & Coerência entre exames
  consecutivos \\
\texttt{websearch} & Busca evidências científicas
  em tempo real & Validar diretrizes clínicas
  do SUS \\
\texttt{sqlite} & Armazena metadados dos pipelines
  e resultados históricos & Histórico de execuções
  TDD \\
\texttt{github} & Versiona código, gerencia
  pull requests e CI/CD & Gatilho de validação
  pré-commit \\
\texttt{pdf} & Gera relatórios de validação
  e conformidade & Relatório SaMD
  para ANVISA \\
\texttt{scihub} & Acesso a artigos científicos
  para ancorar decisões & Citação de
  evidência clínica \\
\midrule
\texttt{code-reviewer} (agente) & Audita qualidade do código
  e detecta vulnerabilidades & Revisar \texttt{pipeline.py}
  antes do merge \\
\texttt{test-engineer} (agente) & Executa suítes TDD e
  gera relatórios de cobertura\footnotemark & Rodar 312 CTs
  do SUS-Twin \\
\bottomrule
\end{tabular}
\end{table}
\footnotetext[1]{Ambiente de execução controlado que isola o código do sistema
operacional hospitalar, prevenindo que falhas em scripts de validação
comprometam dados de pacientes.}
\footnotetext[2]{Cobertura de testes é a métrica que indica qual porcentagem do
código-fonte foi exercitada pelos testes automatizados. O SUS-Twin exige
cobertura mínima de 90\% para conformidade SaMD.}

\subsection{Registro de um novo agente MCP}

Para adicionar um servidor MCP personalizado ao ecossistema, utiliza-se o
comando \texttt{opencode mcp register} seguido de um manifesto JSON:

\begin{lstlisting}[language=bash, caption={Registro de um novo MCP para validação
de segmentação CBCT}, label={lst:mcp-register}, style=bashstyle]
opencode mcp register --manifest cbct-validator.json

# Conteudo de cbct-validator.json:
# {
#   "name": "cbct-validator",
#   "version": "1.0.0",
#   "tools": [
#     {
#       "name": "validate_dicom",
#       "description": "Valida metadados DICOM de tomografia CBCT",
#       "input": { "filepath": "string" },
#       "output": { "is_valid": "boolean", "errors": "string[]" }
#     }
#   ],
#   "transport": "stdio"
# }
\end{lstlisting}

Após o registro, o novo MCP aparece na tabela de roteamento do
OrchestrationEngine\footnote{Motor central de distribuição de tarefas que
escolhe o agente mais adequado para cada tipo de processamento e equilibra
a carga entre todos os servidores MCP disponíveis.} e pode ser invocado por
qualquer skill ou pipeline.


% =============================================================================
\section{Orquestração de Agentes de Validação}
\label{sec:opencode-orquestracao}

A orquestração\footnote{Coordenação automatizada de múltiplos agentes de
software que trabalham em conjunto, como um maestro regendo uma orquestra:
cada agente executa sua parte e o ecossistema garante que todos estejam
sincronizados.} de validação no OpenCode segue um pipeline de três estágios
consecutivos.

\subsection{Estágio 1: Autoevolve — diagnóstico contínuo}

O agente \texttt{autoevolve} varre o ecossistema diariamente em busca de
lacunas de conhecimento ou desempenho. Ele executa o pipeline
\texttt{PLAN~$\rightarrow$~ACT~$\rightarrow$~REFLECT~$\rightarrow$~EXTRACT~$\rightarrow$~EVOLVE}:

\begin{lstlisting}[language=bash, caption={Execução do agente Autoevolve para
diagnóstico do pipeline SUS-Twin}, label={lst:autoevolve}, style=bashstyle]
# Executar autoevolve manualmente
opencode agent run autoevolve --target sus-twin

# Saida esperada:
# [AUTOVOLVE] Iniciando diagnostico do escopo 'sus-twin'...
# [PLAN]     Identificadas 3 lacunas de validacao
# [ACT]      Gerando 2 novas skills complementares
# [REFLECT]  Score de cobertura atual: 87% -> 94%
# [EXTRACT]  Insights persistidos em evolution/insights.json
# [EVOLVE]   Ecossistema atualizado com sucesso
\end{lstlisting}

\subsection{Estágio 2: Code Reviewer — auditoria estática}

O agente \texttt{code-reviewer} inspeciona o código-fonte contra as regras
definidas no SDD (System Design Document)\footnote{Documento de Especificação
de Sistema que formaliza a arquitetura matemática do software, definindo
entradas, saídas, pré-condições, pós-condições e invariantes que o código
nunca pode violar.}:

\begin{lstlisting}[language=bash, caption={Auditoria de código com
Code Reviewer}, label={lst:code-reviewer}, style=bashstyle]
# Revisar todo o diretorio de pipelines
opencode agent run code-reviewer --path ./pipelines/ --strictness high

# Saida resumida:
# [CODE-REVIEW] Auditando 12 arquivos...
# [PASS]        pipelines/segmentacao.py — invariante SDD-I3 respeitada
# [FAIL]        pipelines/preditivo.py — linha 47: tipo inconsistente
# [PASS]        pipelines/telemetria.py — cobertura >= 90%
# Resumo: 10/12 aprovados (83%)

# Corrigir automaticamente falhas simples
opencode agent run code-reviewer --fix --path ./pipelines/preditivo.py
\end{lstlisting}

\subsection{Estágio 3: Test Engineer — validação dinâmica}

O \texttt{test-engineer} executa as suítes TDD completas e gera relatórios
estruturados:\footnote{O relatório inclui tempo de execução, cobertura de
código, falhas detectadas e sugestões de correção, tudo exportável para
formatos JSON e HTML.}

\begin{lstlisting}[language=bash, caption={Execução da suíte TDD completa com
Test Engineer}, label={lst:test-engineer}, style=bashstyle]
# Executar todas as suítes TDD do SUS-Twin
opencode agent run test-engineer --suite sus-twin --report json

# Saida:
# [TEST-ENGINEER] Carregando 15 suites (312 CTs)...
# [RUN]   Suite 01 - OrchestrationEngine (21 CTs) ... 21/21 PASS
# [RUN]   Suite 02 - MultiReasoningEngine (12 CTs) ... 12/12 PASS
# [RUN]   Suite 03 - ScientificProductionAgent (10 CTs) ... 10/10 PASS
# [RUN]   Suite 04 - SROIScanner (8 CTs) ... 8/8 PASS
# [RUN]   Suite 05 - SaMD Compliance (10 CTs) ... 10/10 PASS
# [RUN]   Suite 06 - Scanners Pipeline (10 CTs) ... 10/10 PASS
# [RUN]   Suite 07-15 - Demais modulos ... 241/241 PASS
# ============================================================
# RESUMO: 312/312 PASS   |   Status: APROVADO
# Cobertura: 94.2%       |   Tempo: 14.3s
# ============================================================
\end{lstlisting}


% =============================================================================
\section{Execução de TDD/SDD no Pipeline SUS-Twin}
\label{sec:opencode-tdd}

O Framework SUS-Twin adota uma disciplina binária de engenharia: todo
componente deve ser precedido por uma \textbf{SPEC} (Specification-Driven
Development) e validado por \textbf{testes automatizados} (Test-Driven
Development).\footnote{Essa combinação garante que cada linha de código tenha
uma especificação formal e um teste que prove seu correto funcionamento,
característica essencial para softwares classificados como dispositivos
médicos (SaMD).}

\subsection{O validador central: \texttt{tdd\_validate.py}}

O ecossistema disponibiliza o script \texttt{tdd\_validate.py} como ponto
único de entrada para validação. O \autoref{lst:tdd-validate} mostra a
implementação real que você deve usar em seus projetos:

\begin{lstlisting}[language=Python, caption={Validador TDD central do
Framework SUS-Twin}, label={lst:tdd-validate}]
#!/usr/bin/env python3
"""
tdd_validate.py — Validador TDD central do Framework SUS-Twin.

Uso:
    python tdd_validate.py --suite <nome>      # Executa uma suite especifica
    python tdd_validate.py --all               # Executa todas as suites
    python tdd_validate.py --report json       # Relatorio em JSON
    python tdd_validate.py --chapter 9         # Valida este capitulo

Requer: Python 3.10+, opencode-cli instalado.
"""
import argparse
import json
import subprocess
import sys
from pathlib import Path
from datetime import datetime
from typing import Dict, List, Tuple


# Mapeamento de suites TDD do SUS-Twin
SUITES: Dict[str, Dict] = {
    "orchestration": {
        "path": "tdd/orchestration_engine/",
        "tests": 21,
        "description": "OrchestrationEngine (QE)"
    },
    "reasoning": {
        "path": "tdd/multi_reasoning/",
        "tests": 12,
        "description": "MultiReasoningEngine (MRE)"
    },
    "scientific": {
        "path": "tdd/scientific_production/",
        "tests": 10,
        "description": "ScientificProductionAgent (SPA)"
    },
    "sroi": {
        "path": "tdd/sroi_scanner/",
        "tests": 8,
        "description": "SROIScanner (SS)"
    },
    "samd": {
        "path": "tdd/samd_compliance/",
        "tests": 10,
        "description": "SaMD Compliance Engine"
    },
    "scanners": {
        "path": "tdd/scanners_pipeline/",
        "tests": 10,
        "description": "Scanners Epistemologicos Pipeline"
    },
}


def discover_suites(base_path: Path) -> List[Dict]:
    """Descobre suites TDD disponiveis no diretorio."""
    available = []
    for name, meta in SUITES.items():
        suite_dir = base_path / meta["path"]
        if suite_dir.exists():
            available.append({"name": name, **meta})
    return available


def run_suite(suite_dir: Path, suite_name: str) -> Tuple[int, int, str, str]:
    """
    Executa uma suite TDD via pytest e retorna:
    (passed, total, status, output).

    Utiliza subprocess para garantir isolamento entre suites.
    """
    result = subprocess.run(
        [sys.executable, "-m", "pytest", str(suite_dir),
         "-v", "--tb=short", "--no-header"],
        capture_output=True, text=True, timeout=120
    )
    output = result.stdout + result.stderr

    # Parse do resultado
    passed = output.count("PASSED")
    failed = output.count("FAILED") + output.count("ERROR")
    total = passed + failed

    status = "PASS" if failed == 0 else "FAIL"
    return passed, total, status, output


def generate_report(results: List[Dict], output_format: str = "text"):
    """Gera relatorio formatado (texto ou JSON)."""
    total_passed = sum(r["passed"] for r in results)
    total_tests = sum(r["total"] for r in results)
    timestamp = datetime.now().isoformat()

    if output_format == "json":
        report = {
            "timestamp": timestamp,
            "framework": "SUS-Twin TDD Validator v2.0",
            "summary": {
                "total_tests": total_tests,
                "passed": total_passed,
                "failed": total_tests - total_passed,
                "status": "APPROVED" if total_passed == total_tests else "REJECTED"
            },
            "suites": results
        }
        return json.dumps(report, indent=2)

    # Formato texto (padrao)
    lines = [
        "=" * 60,
        "  SUS-Twin TDD Validation Report",
        f"  {timestamp}",
        "=" * 60,
        ""
    ]
    for r in results:
        status_icon = "[PASS]" if r["status"] == "PASS" else "[FAIL]"
        lines.append(
            f"  {status_icon} {r['description']}: "
            f"{r['passed']}/{r['total']}"
        )
    lines.extend([
        "",
        "-" * 60,
        f"  Total: {total_passed}/{total_tests}",
        f"  Status: {'APROVADO' if total_passed == total_tests else 'REJEITADO'}",
        "-" * 60
    ])
    return "\n".join(lines)


def main():
    parser = argparse.ArgumentParser(
        description="Validador TDD do Framework SUS-Twin"
    )
    parser.add_argument(
        "--suite", type=str,
        help="Nome da suite (orchestration, reasoning, etc.)"
    )
    parser.add_argument("--all", action="store_true",
                        help="Executa todas as suites")
    parser.add_argument("--report", choices=["text", "json"],
                        default="text", help="Formato do relatorio")
    parser.add_argument("--chapter", type=int,
                        help="Valida o conteudo do capitulo especificado")
    args = parser.parse_args()

    base_path = Path(__file__).parent / "tdd"
    if not base_path.exists():
        print("[ERRO] Diretorio 'tdd/' nao encontrado.")
        print("Execute 'opencode init --ecosystem validation' para criar.")
        sys.exit(1)

    available = discover_suites(base_path)

    if args.all or (not args.suite and not args.chapter):
        suites_to_run = available
    elif args.suite:
        suites_to_run = [s for s in available if s["name"] == args.suite]
        if not suites_to_run:
            print(f"[ERRO] Suite '{args.suite}' nao encontrada.")
            sys.exit(1)
    elif args.chapter:
        # Mapeamento: capitulo 9 -> suites de orquestracao
        chapter_map = {9: ["orchestration", "scanners", "samd"]}
        suite_names = chapter_map.get(args.chapter, list(SUITES.keys()))
        suites_to_run = [s for s in available if s["name"] in suite_names]

    results = []
    for suite in suites_to_run:
        print(f"\n>>> Executando: {suite['description']}...")
        passed, total, status, output = run_suite(
            base_path / suite["path"], suite["name"]
        )
        results.append({
            "name": suite["name"],
            "description": suite["description"],
            "passed": passed,
            "total": total,
            "status": status
        })
        if args.report == "text":
            print(output[-500:] if len(output) > 500 else output)

    report = generate_report(results, args.report)
    print(f"\n{report}")

    # Codigo de saida: 0 se aprovado, 1 se rejeitado
    all_passed = all(r["status"] == "PASS" for r in results)
    sys.exit(0 if all_passed else 1)


if __name__ == "__main__":
    main()
\end{lstlisting}

\subsection{Interpretação da saída e correção de falhas}

Ao executar o validador, a saída segue o formato de relatório do
\autoref{lst:tdd-output-example}. Cada falha deve ser tratada com o ciclo
\textbf{Red-Green-Refactor}\footnote{Disciplina central do TDD: (1) escreva
um teste que falha --- Red, (2) implemente o mínimo de código para passar ---
Green, (3) melhore a qualidade sem quebrar o teste --- Refactor.}:

\begin{lstlisting}[style=bashstyle, caption={Exemplo de saída do validador
com falha e correção}, label={lst:tdd-output-example}]
$ python tdd_validate.py --suite orchestration

>>> Executando: OrchestrationEngine (QE)...
  [PASS] TC01: Registers a healthy agent with skills
  [PASS] TC02: Registers multiple agents and validates routing table
  [FAIL] TC03: Unhealthy agents are excluded from routing
  >>> Motivo: agente com health_score=0 ainda é roteado
  >>> Correcao sugerida: adicionar filtro 'if agent.health_score > 0'
  [PASS] TC04: Dispatches task to healthy capable agent
  ...

------------------------------------------------------------
  Total: 20/21
  Status: REJEITADO
------------------------------------------------------------

# Apos aplicar a correcao sugerida:
$ python tdd_validate.py --suite orchestration
  ...
  [PASS] TC03: Unhealthy agents are excluded from routing
  ...
------------------------------------------------------------
  Total: 21/21
  Status: APROVADO
------------------------------------------------------------
\end{lstlisting}

\subsection{Integração SDD: validação de invariantes}

O SDD do SUS-Twin define invariantes\footnote{Regras que devem permanecer
verdadeiras em todos os momentos de execução do sistema. Funcionam como leis
fixas que o software nunca pode violar.} que são verificados em cada execução.
O \autoref{lst:invariants} mostra como validar programaticamente essas regras:

\begin{lstlisting}[language=Python, caption={Validação de invariantes SDD
no pipeline de orquestração}, label={lst:invariants}]
"""
invariances.py — Verificacao de invariantes SDD.

Cada invariante e uma funcao que retorna (nome, resultado, mensagem).
"""
from typing import Callable, List, Tuple

InvariantCheck = Callable[[], Tuple[str, bool, str]]


def check_oe_i1(agents: List[dict]) -> Tuple[str, bool, str]:
    """
    OE-I1: Para toda tarefa em execucao, a carga do agente
    nao pode exceder sua capacidade maxima.
    """
    for agent in agents:
        if agent["current_load"] > agent["max_capacity"]:
            return (
                "OE-I1",
                False,
                f"Agente {agent['id']}: carga {agent['current_load']} > "
                f"capacidade {agent['max_capacity']}"
            )
    return ("OE-I1", True, "Todos os agentes dentro do limite de carga")


def check_mre_i1(context: dict) -> Tuple[str, bool, str]:
    """
    MRE-I1: O nivel de confianca do raciocinio deve ser >= 0.6
    para ativar modo dedutivo.
    """
    if context.get("confidence", 0) < 0.6:
        return (
            "MRE-I1",
            False,
            f"Confianca {context.get('confidence', 0)} abaixo do limiar 0.6"
        )
    return ("MRE-I1", True, "Confianca adequada para modo dedutivo")


def validate_all_invariants(agents, context) -> List[Tuple[str, bool, str]]:
    """Executa todas as verificacoes de invariante."""
    checks = [
        check_oe_i1(agents),
        check_mre_i1(context),
    ]
    return checks


if __name__ == "__main__":
    # Simulacao de agentes e contexto
    mock_agents = [
        {"id": "agent-01", "current_load": 3, "max_capacity": 5},
        {"id": "agent-02", "current_load": 6, "max_capacity": 5},  # violacao
    ]
    mock_context = {"confidence": 0.8}

    results = validate_all_invariants(mock_agents, mock_context)
    all_pass = True
    for name, passed, msg in results:
        status = "[PASS]" if passed else "[FAIL]"
        print(f"  {status} {name}: {msg}")
        if not passed:
            all_pass = False

    print(f"\nInvariantes: {'TODAS OK' if all_pass else 'FALHAS DETECTADAS'}")
    # Saida:
    #   [PASS] OE-I1: Agente agent-01 dentro do limite
    #   [FAIL] OE-I1: Agente agent-02: carga 6 > capacidade 5
    #   [PASS] MRE-I1: Confianca adequada para modo dedutivo
    #   Invariantes: FALHAS DETECTADAS
\end{lstlisting}


% =============================================================================
\section{Pipeline CI/CD para Validação Contínua}
\label{sec:opencode-cicd}

A validação não deve ser um evento pontual, mas um processo contínuo
integrado ao fluxo de desenvolvimento. O OpenCode oferece integração nativa
com Git hooks\footnote{Scripts que o Git executa automaticamente em eventos
como \texttt{pre-commit} (antes de um commit) ou \texttt{pre-push} (antes de
um push), permitindo barrar código que não passe nos testes.} e GitHub
Actions.

\subsection{Git hook de pré-commit}

O hook \texttt{pre-commit} abaixo executa o validador TDD antes de cada
commit, bloqueando alterações que quebrem os testes:

\begin{lstlisting}[language=bash, caption={Git hook pre-commit para validacao
automatica}, label={lst:pre-commit}, style=bashstyle]
#!/bin/bash
# .git/hooks/pre-commit — Validacao pre-commit do SUS-Twin

echo "=== SUS-Twin: Validacao pre-commit ==="

# Executa apenas as suites afetadas pelos arquivos modificados
CHANGED=$(git diff --cached --name-only --diff-filter=ACM)

if echo "$CHANGED" | grep -q "tdd/"; then
    echo "Alteracoes detectadas em suites TDD. Executando validacao completa..."
    python tdd_validate.py --all --report text
elif [ -n "$CHANGED" ]; then
    echo "Alteracoes em codigo fonte. Executando suites relacionadas..."
    python tdd_validate.py --all --report text
else
    echo "Nenhuma alteracao relevante. Pulando validacao."
    exit 0
fi

# Se o script retornar 0, o commit prossegue; se retornar 1, e bloqueado
if [ $? -ne 0 ]; then
    echo "=== VALIDACAO FALHOU: Commit BLOQUEADO ==="
    echo "Corrija as falhas antes de commitar."
    exit 1
fi

echo "=== VALIDACAO OK: Commit permitido ==="
exit 0
\end{lstlisting}

\subsection{Integração com GitHub Actions}

Para ambientes colaborativos, o pipeline CI/CD deve ser definido em
\texttt{.github/workflows/validate.yml}:

\begin{lstlisting}[basicstyle=\ttfamily\scriptsize, caption={Workflow GitHub Actions para
validacao continua do SUS-Twin}, label={lst:github-actions}]
name: SUS-Twin Validation Pipeline

on:
  push:
    branches: [ main, develop ]
  pull_request:
    branches: [ main ]

jobs:
  validate:
    runs-on: ubuntu-latest
    steps:
      - uses: actions/checkout@v4

      - name: Setup Python
        uses: actions/setup-python@v5
        with:
          python-version: '3.12'

      - name: Install OpenCode CLI
        run: npm install -g @opencode/cli

      - name: Initialize Ecosystem
        run: opencode init --ecosystem validation

      - name: Run TDD Validation
        id: tdd
        run: |
          python tdd_validate.py --all --report json > report.json
          echo "status=$(python -c \
            'import json; r=json.load(open("report.json"));
             print(r["summary"]["status"])')" >> $GITHUB_OUTPUT

      - name: Generate Validation Badge
        if: always()
        run: |
          python -c "
          import json
          r = json.load(open('report.json'))
          status = r['summary']['status']
          passed = r['summary']['passed']
          total = r['summary']['total']
          color = 'brightgreen' if status == 'APPROVED' else 'red'
          url = f'https://img.shields.io/badge/TDD-{passed}%2F{total}-{color}'
          print(f'SUS-Twin TDD: {passed}/{total}')
          with open('badge_url.txt', 'w') as f:
              f.write(url)
          "

      - name: Upload Report
        if: always()
        uses: actions/upload-artifact@v4
        with:
          name: tdd-report
          path: report.json
\end{lstlisting}


% =============================================================================
\section{Badge de Validação e Relatórios}
\label{sec:opencode-badge}

Ao final de cada pipeline de validação, o ecossistema gera um badge
visual\footnote{Distintivo digital que exibe o status atual dos testes,
similar aos badges de CI/CD (como \texttt{build passing}), usado em
repositórios GitHub para comunicar rapidamente a saúde do projeto.} que
pode ser inserido em README ou relatórios. O \autoref{lst:badge-gen} mostra
como gerar programaticamente:

\begin{lstlisting}[language=Python, caption={Gerador de badge de validacao
para o Framework SUS-Twin}, label={lst:badge-gen}]
#!/usr/bin/env python3
"""
badge_gen.py — Gerador de badge de validacao para SUS-Twin.

Gera um badge SVG local e uma URL shields.io para README.
"""
import json
from pathlib import Path
from datetime import datetime


def generate_badge(passed: int, total: int, output_path: Path):
    """Gera badge de validacao no formato SVG."""
    ratio = passed / total if total > 0 else 0
    color = "brightgreen" if ratio == 1.0 else \
            "yellow" if ratio >= 0.9 else "red"

    label = f"SUS-Twin TDD"
    value = f"{passed}/{total}"
    svg = f'''<?xml version="1.0" encoding="UTF-8"?>
<svg xmlns="http://www.w3.org/2000/svg" width="180" height="20">
  <linearGradient id="b" x2="0" y2="100%">
    <stop offset="0" stop-color="#bbb" stop-opacity=".1"/>
    <stop offset="1" stop-opacity=".1"/>
  </linearGradient>
  <rect rx="3" width="180" height="20" fill="#555"/>
  <rect rx="3" x="85" width="95" height="20" fill="{color}"/>
  <g fill="#fff" font-family="DejaVu Sans" font-size="11">
    <text x="6" y="15">{label}</text>
    <text x="95" y="15">{value}</text>
  </g>
</svg>'''
    output_path.write_text(svg, encoding="utf-8")
    print(f"[OK] Badge salvo em: {output_path}")


def get_shields_url(passed: int, total: int) -> str:
    """Retorna URL do badge via shields.io."""
    ratio = passed / total if total > 0 else 0
    color = "brightgreen" if ratio == 1.0 else \
            "yellow" if ratio >= 0.9 else "red"
    label = "SUS-Twin%20TDD"
    value = f"{passed}%2F{total}"
    return (f"https://img.shields.io/badge/{label}-{value}-{color}")


if __name__ == "__main__":
    report_path = Path("report.json")
    if not report_path.exists():
        print("[ERRO] report.json nao encontrado. Execute tdd_validate.py primeiro.")
        exit(1)

    report = json.loads(report_path.read_text(encoding="utf-8"))
    summary = report["summary"]

    # Gera badge SVG local
    generate_badge(summary["passed"], summary["total"],
                   Path("badge_sustwin.svg"))

    # Gera URL shields.io para README
    url = get_shields_url(summary["passed"], summary["total"])
    print(f"[INFO] URL para README.md:")
    print(f"  ![TDD]({url})")
\end{lstlisting}

O badge gerado pode ser inserido no README do repositório:

\begin{lstlisting}[basicstyle=\ttfamily\scriptsize, caption={Inserção do badge no README.md},
label={lst:badge-readme}]
# Framework SUS-Twin

![SUS-Twin TDD](https://img.shields.io/badge/SUS--Twin%20TDD-312%2F312-brightgreen)

Pipeline de Gêmeos Digitais Periodontais validado pelo OpenCode Ecosystem.
\end{lstlisting}

\subsection{Badge ASCII do capítulo}

Como síntese do que foi aprendido, o badge ASCII abaixo resume as
competências adquiridas neste capítulo:

\begin{highlight}
\begin{verbatim}
+----------------------------------------------------------+
| Badge do Capítulo                                          |
+----------------------------------------------------------+
| Nivel-alvo: N2 — Pesquisa Reprodutivel                    |
| Habilidades: orquestrar agentes, executar TDD,            |
|              configurar pipeline CI/CD                    |
| Pre-requisitos: Python 3.11+, OpenCode CLI, Git           |
| Comando de verificacao:                                   |
| tdd_validate.py --chapter 9                                |
+----------------------------------------------------------+
\end{verbatim}
\end{highlight}


% =============================================================================
\section{Exercícios Práticos}
\label{sec:opencode-exercicios}

Os exercícios a seguir consolidam o aprendizado prático do capítulo. Cada um
deve ser resolvido em um terminal real com o OpenCode Ecosystem instalado.

\subsection*{Exercício 1 — Configurar um agente MCP personalizado}

\textbf{Objetivo:} Registrar um novo servidor MCP para validação de arquivos
DICOM odontológicos.

\begin{enumerate}
  \item Crie o arquivo \texttt{dicom-mcp.json} com o manifesto de um MCP que
        exponha a ferramenta \texttt{validate\_dicom} com parâmetro de entrada
        \texttt{filepath} (string) e saída \texttt{is\_valid} (booleano).
  \item Registre o MCP usando \texttt{opencode mcp register --manifest
        dicom-mcp.json}.
  \item Verifique se o MCP aparece na lista ativa com
        \texttt{opencode mcp list}.
  \item \textbf{Questão:} Qual a importância de ter um MCP dedicado para
        validação DICOM no pipeline SUS-Twin? (\textit{Dica: relacione com
        conformidade SaMD e RDC 657/2022 da ANVISA}).
\end{enumerate}

\subsection*{Exercício 2 — Executar TDD e corrigir uma falha}

\textbf{Objetivo:} Simular um pipeline TDD com falha induzida e corrigi-la.

\begin{enumerate}
  \item Crie um diretório \texttt{tdd/pratica/} com dois arquivos:
        \texttt{test\_exemplo.py} contendo um teste que espera que a função
        \texttt{somar(2, 2)} retorne 4, e \texttt{exemplo.py} contendo uma
        implementação de \texttt{somar} que retorna 5 (falha induzida).
  \item Execute \texttt{python -m pytest tdd/pratica/ -v}. Observe a falha.
  \item Corrija a função \texttt{somar} para retornar o valor correto.
  \item Execute novamente e confirme que o teste passa.
  \item \textbf{Questão:} Explique como o ciclo Red-Green-Refactor se aplica
        a este exercício e por que ele é crucial para software de grau médico.
\end{enumerate}

\subsection*{Exercício 3 — Criar um badge de validação}

\textbf{Objetivo:} Gerar um badge de validação a partir de um relatório TDD.

\begin{enumerate}
  \item Execute o script \autoref{lst:tdd-validate} com \texttt{--report json}
        (ou crie um \texttt{report.json} manual com 5 suítes fictícias, todas
        com 100\% de aprovação).
  \item Execute o script \autoref{lst:badge-gen} para gerar o badge SVG.
  \item Abra o arquivo \texttt{badge\_sustwin.svg} em um navegador para
        confirmar a aparência.
  \item \textbf{Questão:} Quais informações um badge de validação transmite
        para a equipe de desenvolvimento e para órgãos reguladores como a
        ANVISA? Por que o badge deve refletir dados em tempo real e não
        estáticos?
\end{enumerate}

\subsection*{Exercício 4 (Desafio) — Pipeline CI/CD completo}

\textbf{Objetivo:} Integrar todos os conceitos em um pipeline funcional.

\begin{enumerate}
  \item Crie um repositório Git local.
  \item Configure o hook \texttt{pre-commit} do \autoref{lst:pre-commit}.
  \item Crie uma suíte TDD com 3 testes (pode ser simples, como validação
        de soma de números).
  \item Faça um commit que passe nos testes e outro que falhe.
  \item Observe o comportamento do hook em cada caso.
  \item \textbf{Questão:} Por que a validação pré-commit é preferível à
        validação apenas no CI/CD remoto? Cite ao menos dois motivos
        relacionados à produtividade da equipe.
\end{enumerate}

% =============================================================================
% FIM DO CAPÍTULO 9
% =============================================================================
